  % !TEX TS-program = xelatex
  % !TEX encoding = UTF-8 Unicode
  % !Mode:: "TeX:UTF-8"

  \documentclass{resume}
  \usepackage{zh_CN-Adobefonts_external} % Simplified Chinese Support using external fonts (./fonts/zh_CN-Adobe/)
  %\usepackage{zh_CN-Adobefonts_internal} % Simplified Chinese Support using system fonts
  \usepackage{linespacing_fix} % disable extra space before next section
  \usepackage{cite}

  \begin{document}
  \pagenumbering{gobble} % suppress displaying page number

  \name{柳振跃}

  \basicInfo{
    邮箱: jarvisliu1014@gmail.com \textperiodcentered\ 电话/微信: 193-3400-8868 
  }
  \basicInfo{
    GitHub: https://github.com/Jaden-47 \textperiodcentered\ 个人主页: https://jaden-47.github.io/
  }
  
  \section{\faGraduationCap\  教育背景及证书}
  \datedsubsection{\textbf{系统架构设计师}\ 国家计算机技术与软件（高级）专业技术资格}{2024}
  \datedsubsection{\textbf{加州大学圣克鲁兹分校（USNEWS TOP100）}, 计算机科学专业\ 学士}{2019 -- 2022}


  \section{\faUsers\ 工作经历}
  \datedsubsection{\textbf{小米科技有限公司}\ 南京}{2025年4月 -- 至今}
  \role{C/C++, Rust, Flutter, Perfetto}{\textbf{HyperOS 低延迟实时系统框架} 核心调度与性能优化}
  \begin{onehalfspacing}
  系统开发工程师
  \begin{itemize}
    \item 参与 HyperOS 低延迟实时系统框架（RustUI）架构设计与落地，基于 C/C++ 及 Rust FFI 重构核心调度层与事件处理链路，设计跨语言 FFI 边界的低开销数据传输与线程安全机制。
    \item 设计多优先级异步任务调度器，打通 Rust async runtime 与 Flutter engine 的事件循环与线程模型，解耦关键路径与后台任务，消除 head-of-line blocking，P99 尾延迟显著降低，保障系统级场景的确定性低延迟响应。
    \item 优化关键路径延迟：通过 mmap + 异步 IO 实现资源零拷贝加载，设计多优先级任务调度与流控机制（优先级队列 + 动态预算分配），消除冗余计算。端到端延迟降低 70\%+，CPU 开销降低约 20\%。
    \item 基于 LangChain 搭建 AI Agent 驱动的系统性能可观测平台：使用 Perfetto / atrace 插桩关键路径，持续采集微秒级延迟与内存开销数据，自动化输出归因分析与异常预警，前置暴露性能风险。
  \end{itemize}
  \end{onehalfspacing}


  \datedsubsection{\textbf{阿里斑马智行网络有限公司}\ 杭州}{2022年11月 -- 2025年1月}
  \role{C/C++, Rust, eBPF}{\textbf{AliOS DDS} 车载实时通信中间件}
  \begin{onehalfspacing}
  系统中间件开发工程师
  \begin{itemize}
    \item 主导 Discovery 发现模块性能优化：小规模场景发现时间从 130μs 降至 \textasciitilde40μs，大规模场景从 3\textasciitilde5s 降至 300ms，并基于 tokio 重构异步架构，消除长时间持锁导致的尾延迟。
    \item 设计低延迟、确定性的数据传输链路，基于 POSIX shm / mmap 实现共享内存零拷贝通道，配合用户态序列化优化，单条消息端到端延迟控制在 \textasciitilde50μs。
    \item 基于 AF\_XDP socket + eBPF 构建用户态网络栈，绕过内核协议栈实现零拷贝收发包，显著降低上下文切换与数据拷贝开销，为高吞吐低尾延迟场景打下基础。
    \item 参与内核模块适配与实时调度策略调优，确保 DDS 关键链路在 RT 场景下的确定性与低抖动表现。
    \item 基于 Vue 搭建 DDS 中间件可视化管理与配置平台，实现服务拓扑管理、配置热更新与实时监控面板，提升中间件运维效率与易用性。
  \end{itemize}
  \end{onehalfspacing}

  % \datedsubsection{\textbf{分布式科学上网姿势}}{2014年6月 -- 至今}
  % \role{Golang, Linux}{个人项目，和富帅糕合作开发}
  % \begin{onehalfspacing}
  % 分布式负载均衡科学上网姿势, https://github.com/cyfdecyf/cow
  % \begin{itemize}
  %   \item 修复了连接未正常关闭导致文件描述符耗尽的 bug
  %   \item 使用Chord 哈希 URL, 实现稳定可靠地分流
  %   \item xxx (尽量使用量化的客观结果)
  % \end{itemize}
  % \end{onehalfspacing}

  % \datedsubsection{\textbf{\LaTeX\ 简历模板}}{2015 年5月 -- 至今}
  % \role{\LaTeX, Python}{个人项目}
  % \begin{onehalfspacing}
  % 优雅的 \LaTeX\ 简历模板, https://github.com/billryan/resume
  % \begin{itemize}
  %   \item 容易定制和扩展
  %   \item 完善的 Unicode 字体支持，使用 \XeLaTeX\ 编译
  %   \item 支持 FontAwesome 4.5.0
  % \end{itemize}
  % \end{onehalfspacing}

  % Reference Test
  %\datedsubsection{\textbf{Paper Title\cite{zaharia2012resilient}}}{May. 2015}
  %An xxx optimized for xxx\cite{verma2015large}
  %\begin{itemize}
  %  \item main contribution
  %\end{itemize}

  \section{\faCode\ 项目经历}
\datedsubsection{\textbf{matchX}\ 个人项目}{}
\role{Rust, no\_std}{\textbf{低延迟交易匹配引擎}}
\begin{onehalfspacing}
\begin{itemize}
  \item 实现亚微秒级 P99 延迟的订单匹配引擎，Rust no\_std 核心，支持限价/市价/Iceberg/Stop-Limit 等订单类型及 Self-Trade Prevention 机制。
  \item 设计 cache-friendly 数据结构：Order 精确 64B 对齐单一缓存行，Arena 预分配 + 空闲链表复用实现零分配热路径，混合稠密/稀疏订单簿基于 occupancy bitset 实现 O(1) BBO 查询。
  \item 二进制 WAL 日志 + CRC32 校验，异步后台写入；fat LTO + panic=abort + PGO 构建优化；Criterion + HDRHistogram 延迟基准测试体系。
\end{itemize}
\end{onehalfspacing}

\section{\faCogs\ 技能清单}
  % increase linespacing [parsep=0.5ex]
  \begin{itemize}[parsep=0.5ex]
      \item \textbf{语言}：C/C++（精通）, Rust（含 unsafe / async / FFI）, Python
      \item \textbf{前端 \& 框架}：Flutter / RustUI 声明式框架, Vue 前端工程化, 图形渲染管线 (Skia / SurfaceFlinger)
      \item \textbf{系统编程}：Linux 内存管理 (mmap / shm), 异步 IO, CPU 调度与线程亲和性, 内核模块
      \item \textbf{性能分析}：Perfetto, atrace, ftrace, perf, eBPF, gdb
      \item \textbf{内核与网络}：eBPF / AF\_XDP, 用户态协议栈, Zero-Copy, 实时调度 (SCHED\_FIFO / SCHED\_RR)
  \end{itemize}

  % \section{\faHeartO\ 获奖情况}
  % \datedline{\textit{第一名}, xxx 比赛}{2013 年6 月}
  % \datedline{其他奖项}{2015}

  % \section{\faInfo\ 其他}
  % % increase linespacing [parsep=0.5ex]
  % \begin{itemize}[parsep=0.5ex]
  %   \item 技术博客: http://blog.yours.me
  %   \item GitHub: https://github.com/username
  %   \item 语言: 英语 - 熟练(TOEFL xxx)
  % \end{itemize}

  %% Reference
  %\newpage
  %\bibliographystyle{IEEETran}
  %\bibliography{mycite}
  \end{document}
